\section{Workload Corpus and Characterization}
\label{sec:workloads}

\subsection{Sources}
\label{sec:sources}

LSSP's primary Pilot-V2 corpus draws from three independently released
workload sources: BurstGPT \citep{burstgpt2025}, Microsoft's official Azure
2024 LLM-inference trace \citep{azure2024trace,stojkovic2025dynamollm}, and
Alibaba Cloud's Bailian/Qwen anonymized usage-trace
release \citep{bailianqwen2026}. Azure's 2023 conversation/code traces
\citep{patel2024splitwise} and Stage-0's historical windows contribute
additional context but are not part of the primary Pilot-V2 120-window
matrix beyond the 10-per-source Stage-0-reused windows described below.
TraceLab \citep{tracelab2026} is retained as an out-of-distribution,
coding-agent workload family for future expansion, not as one of the three
primary sources, because its independence from a prior 512-window public sweep cannot be fully verified from public artifacts alone, as detailed in the project's provenance records.

As established in \S\ref{sec:related-position}, BurstGPT's own reported
collection methodology logs traffic from a regional customer of Microsoft's
Azure OpenAI GPT service, and this project's Azure trace is Microsoft's own
official release of internal Azure LLM-inference-service traffic. Both
therefore sit on the same underlying cloud provider's infrastructure. We
describe BurstGPT and Azure as distinct, independently released workload
sources -- different collection methodology, different customer
population, different release provenance and version history -- and we do
not describe them as independent providers. Bailian/Qwen (Alibaba Cloud)
is a genuinely different underlying platform from both.

\subsection{Representation and Provenance}
\label{sec:provenance}

Each source is ingested through a dedicated, field-level-audited adapter that maps the source's native
schema onto a common internal request representation (arrival time, input
tokens, output tokens, total tokens, session identifier where available,
model class). No field is fabricated where a source's native schema lacks
it (for example, BurstGPT's documented schema has no tenant, KV-reuse, SLO,
or failure-code field, and none is synthesized). License status for each source was verified before any confirmatory use of a real, full-scale sample (for example, BurstGPT is under CC-BY-4.0).

\subsection{Window Construction}
\label{sec:windows}

The Pilot-V2 corpus freezes exactly 120 windows of 200 requests each: 40
per source (\texttt{azure\_llm\_2024}, \texttt{bailian\_qwen},
\texttt{burstgpt}), each composed of 10 windows reused verbatim from the
Stage-0 pilot (\S\ref{sec:stage0}) plus 30 new windows drawn by a
deterministic, outcome-blind, source-symmetric pinned-extension sampler
that uses only valid-row index ranges and fixed seeds -- never scheduler
metrics, policy identities, or Stage-0 tie/non-tie outcomes (verified by
structural audit to contain no BurstGPT-outcome-targeting signal). Windows
are assigned to EARLY/MIDDLE/LATE chronology strata by deterministic
relative chronology (60/30/30 windows respectively); Stage-0-reused windows
are all EARLY. The freeze passed a full structural validation gate --
unique window identities, exact per-source and Stage-0/new counts, fixed
\texttt{request\_count = 200}, no prohibited overlap, valid chronology, and
absence of any scheduler/policy/load-calibration outcome field in the
selection rule -- recorded with a canonical content hash (excluding
generation timestamp; full value in Table~\ref{tab:hashes}) before any
Pilot-V2 load calibration or scheduler outcome was generated.

\subsection{Descriptors}
\label{sec:descriptors}

Each window is described by a 28-feature descriptor vector spanning prompt-
and output-token-length statistics (mean, p90, long-prompt fraction at the
512/2048/8192-token thresholds), arrival-timing statistics (mean arrival
rate, inter-arrival percentiles), and derived load-pressure proxies. Three
deterministically redundant pressure-proxy features ($R^2=1.0$ against
other included features) are retained in the underlying dataset but
excluded from the classifier input used in \S\ref{sec:separability}.
Prompt-token figures compose differently across sources for architectural
reasons: TraceLab's and Azure's prompt-token fields are potentially
multi-turn-cumulative context, while BurstGPT's is a single per-request
figure. This is a property of how each source's underlying workload is
architected (a growing agent/conversation session versus a single stateless
API call), not an adapter artifact, and it is preserved as an explicit
caveat on every prompt-length-based claim below.

\subsection{Cross-Source Differences}
\label{sec:cross-source-diffs}

This characterization analysis uses a separate, larger population than the
primary 120-window ranking corpus: $100$ windows (window size 200) drawn
from each of \emph{four} source classes -- \texttt{azure\_llm\_2024},
\texttt{bailian\_qwen}, \texttt{burstgpt}, and \texttt{tracelab}
(\S\ref{sec:sources}) -- for $n=400$ windows total. TraceLab is included
here only as a fourth descriptor class to stress-test the classifier
against a maximally distinct, out-of-distribution workload family; it is
not one of the three primary Pilot-V2 ranking-portability sources and
contributes no windows to the 120-window corpus used in
\S\ref{sec:results}--\S\ref{sec:results-slo}. Under a leave-one-time-bucket-out (EARLY/MIDDLE/LATE), grouped,
chronologically blocked cross-validation over these $n=400$ windows, a
random forest classifier distinguishes the four
descriptor-labeled source classes with pooled balanced accuracy
$1.000$ (mean $\pm$ std across three folds: $1.000 \pm 0.000$); a
multinomial logistic regression reaches $0.988$ balanced accuracy
($0.987 \pm 0.007$ across folds); and a depth-limited (depth $\le 3$)
decision tree reaches $0.980$ balanced accuracy ($0.980 \pm 0.018$ across
folds, worst single fold $0.955$ with LATE held out). This grouped,
chronologically blocked evaluation supersedes an earlier random-split
result (which agreed at balanced accuracy $1.0$) for any manuscript claim,
because random splitting is vulnerable to within-window-adjacent-time leakage that the grouped evaluation is structurally immune to (and was verified to prevent data leakage across the separability pipeline).

\subsection{Source Separability}
\label{sec:separability}

The workload sources are confirmed measurably separable, with the caveats
below. The safe claim
supported by this evidence is that the workload sources' observed
descriptor distributions are measurably different, and that source identity
is strongly separable under a leakage-resistant, chronologically blocked
evaluation. This is not evidence that the underlying providers'
conversations or tasks intrinsically differ in length or structure, only
that their \emph{observed, request-level} descriptor figures do, for a
mixture of architectural and possibly-genuine-content reasons that this
classifier-based analysis cannot separate (\S\ref{sec:descriptors}).

\subsection{Drivers of Differences}
\label{sec:drivers}

An attribution analysis triangulating held-out permutation importance,
single-feature balanced accuracy, leave-one-feature-out deltas, and
feature-group-ablation deltas (all under the same grouped chronological
cross-validation) identifies two tiers of driving features. Tier 1 --
incrementally load-bearing, with nonzero permutation importance and a
measurable accuracy cost when the \texttt{prompt\_length} feature group is
removed ($-0.0051$) -- consists of \texttt{long\_prompt\_fraction\_512},
\texttt{long\_prompt\_fraction\_2048}, and \texttt{prompt\_tokens\_mean}.
Tier 2 -- individually strong single-feature classifiers (balanced accuracy
up to $0.980$) whose removal, individually or as a whole feature group,
costs $0.000$ accuracy because their signal is redundant with Tier 1 --
includes total/prompt-token percentile features, inter-arrival percentiles,
and mean arrival rate. Prompt-length structure, specifically the fraction
of long prompts at the 512- and 2048-token thresholds and mean prompt
length, is therefore the classifier's primary, non-redundant predictive
feature for source separability, in the ablation sense defined above
(largest classifier accuracy cost when removed) -- not
a claim that prompt length is the true underlying cause of the sources'
differences, which may also reflect the multi-turn-cumulative-vs-single-request
token-accounting confound already noted in \S\ref{sec:descriptors};
arrival-timing structure and output length are individually
predictive but carry no additional signal once prompt-length features are
present. A critical-ablation check confirms this separability is broad
rather than dependent on any single feature: removing the two strongest
Tier-1 features together changes random-forest balanced accuracy by at
most $-0.0025$.

\subsection{Window-Size Sensitivity}
\label{sec:window-size}

The Pilot-V2 corpus uses a fixed window size of 200 requests throughout,
matching Stage-0 and the LLM~2026 comparison window size for
methodological consistency. A dedicated window-size sensitivity sweep --
varying the number of requests per window and re-checking
separability/discriminability, distinct from the frozen
\texttt{WINDOW\_SIZE\_SENSITIVITY} robustness-contract item
(\S\ref{sec:results-robustness}), which despite the shared name concerns
how many \emph{windows} are subsampled, not requests per window -- has not
yet been run; this is noted as an open item in \S\ref{sec:threats} rather
than a result reported here.
